\begin{definition}
  \label{def:C}
  \lean{C}
  \leanok
  For any natural number $n \in \mathbb{N}$, we define the map $C: \mathbb{N} \to \mathbb{N}$ as:
  \[ C(n) = \lfloor 3n/2 \rfloor. \]
\end{definition}

\begin{definition}
  \label{def:D}
  \lean{D}
  \leanok
  For any natural number $n \in \mathbb{N}$, we define the map $D: \mathbb{N} \to \mathbb{N}$ as:
  \[ D(n) = \lceil 3n/2 \rceil. \]
\end{definition}

\begin{lemma}
  \label{lemma:C_even}
  \lean{C_even}
  \leanok
  If $n \in \mathbb{N}$ is even, then $C(n) = 3n/2$.
  \begin{proof}
    \leanok
    If $n$ is even, then $n = 2k$ for some $k \in \mathbb{N}$. Then $3n/2 = 3(2k)/2 = 3k$, which is an integer.
    Thus $C(n) = \lfloor 3k \rfloor = 3k = 3n/2$.
  \end{proof}
\end{lemma}

\begin{lemma}
  \label{lemma:C_odd}
  \lean{C_odd}
  \leanok
  If $n \in \mathbb{N}$ is odd, then $C(n) = (3n-1)/2$.
  \begin{proof}
    \leanok
    If $n$ is odd, then $n = 2k+1$ for some $k \in \mathbb{N}$.
    Then $C(n) = \lfloor 3(2k+1)/2 \rfloor = \lfloor (6k+3)/2 \rfloor = \lfloor 3k + 1.5 \rfloor = 3k+1$.
    Also, $(3n-1)/2 = (3(2k+1)-1)/2 = (6k+2)/2 = 3k+1$.
  \end{proof}
\end{lemma}

\begin{lemma}
  \label{lemma:D_even}
  \lean{D_even}
  \leanok
  If $n \in \mathbb{N}$ is even, then $D(n) = 3n/2$.
  \begin{proof}
    \leanok
    If $n$ is even, then $n = 2k$ for some $k \in \mathbb{N}$. Then $3n/2 = 3k$, which is an integer.
    Thus $D(n) = \lceil 3k \rceil = 3k = 3n/2$.
  \end{proof}
\end{lemma}

\begin{lemma}
  \label{lemma:D_odd}
  \lean{D_odd}
  \leanok
  If $n \in \mathbb{N}$ is odd, then $D(n) = (3n+1)/2$.
  \begin{proof}
    \leanok
    If $n$ is odd, then $n = 2k+1$ for some $k \in \mathbb{N}$.
    Then $D(n) = \lceil 3(2k+1)/2 \rceil = \lceil 3k + 1.5 \rceil = 3k+2$.
    Also, $(3n+1)/2 = (3(2k+1)+1)/2 = (6k+4)/2 = 3k+2$.
  \end{proof}
\end{lemma}

\begin{lemma}
  \label{lemma:C_formula}
  \lean{C_formula}
  \leanok
  For all $n \in \mathbb{N}$, the integer value of $C(n)$ is given by:
  \[ C(n) = \frac{6n - 1 + (-1)^n}{4}. \]
  \begin{proof}
    \leanok
    We consider two cases for the parity of $n$.
    If $n$ is even, then $(-1)^n = 1$ and $C(n) = 3n/2$. The formula gives $(6n-1+1)/4 = 6n/4 = 3n/2$.
    If $n$ is odd, then $(-1)^n = -1$ and $C(n) = (3n-1)/2$. The formula gives $(6n-1-1)/4 = (6n-2)/4 = (3n-1)/2$.
  \end{proof}
\end{lemma}

\begin{lemma}
  \label{lemma:D_formula}
  \lean{D_formula}
  \leanok
  For all $n \in \mathbb{N}$, the integer value of $D(n)$ is given by:
  \[ D(n) = \frac{6n + 1 - (-1)^n}{4}. \]
  \begin{proof}
    \leanok
    If $n$ is even, then $(-1)^n = 1$ and $D(n) = 3n/2$. The formula gives $(6n+1-1)/4 = 6n/4 = 3n/2$.
    If $n$ is odd, then $(-1)^n = -1$ and $D(n) = (3n+1)/2$. The formula gives $(6n+1+1)/4 = (6n+2)/4 = (3n+1)/2$.
  \end{proof}
\end{lemma}

\begin{definition}
  \label{def:Citer}
  \lean{Citer}
  \leanok
  For any $i, n \in \mathbb{N}$, we denote the $i$-th iterate of $C$ by $C_i(n) = C^i(n)$.
\end{definition}

\begin{definition}
  \label{def:Diter}
  \lean{Diter}
  \leanok
  For any $i, n \in \mathbb{N}$, we denote the $i$-th iterate of $D$ by $D_i(n) = D^i(n)$.
\end{definition}

\begin{lemma}
  \label{lemma:C_succ_eq_D_succ}
  \lean{C_succ_eq_D_succ}
  \leanok
  For all $m \in \mathbb{N}$, we have $C(m+1) = D(m) + 1$.
  \begin{proof}
    \leanok
    We calculate:
    $C(m+1) = \lfloor \frac{3(m+1)}{2} \rfloor = \lfloor \frac{3m+3}{2} \rfloor = \lfloor \frac{3m+1}{2} + 1 \rfloor = \lfloor \frac{3m+1}{2} \rfloor + 1$.
    Since $D(m) = \lceil \frac{3m}{2} \rceil = \lfloor \frac{3m+1}{2} \rfloor$, we have $C(m+1) = D(m) + 1$.
  \end{proof}
\end{lemma}

\begin{lemma}
  \label{lemma:Citer_succ_eq_Diter_succ}
  \lean{Citer_succ_eq_Diter_succ}
  \leanok
  For all $i, n \in \mathbb{N}$, we have $C_i(n+1) = D_i(n) + 1$.
  \begin{proof}
    \leanok
    We proceed by induction on $i$.
    Base case $i=0$: $C_0(n+1) = n+1$ and $D_0(n)+1 = n+1$.
    Inductive step: Assume $C_i(n+1) = D_i(n) + 1$. Then
    $C_{i+1}(n+1) = C(C_i(n+1)) = C(D_i(n) + 1)$.
    By Lemma~\ref{lemma:C_succ_eq_D_succ}, $C(D_i(n) + 1) = D(D_i(n)) + 1 = D_{i+1}(n) + 1$.
  \end{proof}
\end{lemma}

\begin{definition}
  \label{def:A}
  \lean{A}
  \leanok
  We define the map $A: \mathbb{N} \times \mathbb{N} \to \mathbb{N} \times \mathbb{N}$ as follows:
  Let $(a, b) \in \mathbb{N} \times \mathbb{N}$ and $n = \lfloor a/2 \rfloor$. Then:
  \[
  A(a, b) = \begin{cases}
    (3n+2, b+2) & \text{if } a \text{ is even}, \\
    (3n+3, b-1) & \text{if } a \text{ is odd}.
  \end{cases}
  \]
\end{definition}

\begin{definition}
  \label{def:Aiter}
  \lean{Aiter}
  \leanok
  For any $i \in \mathbb{N}$ and $(a, b) \in \mathbb{N} \times \mathbb{N}$, we denote the $i$-th iterate of $A$ by $A_i(a, b) = A^i(a, b)$.
\end{definition}

\begin{lemma}
  \label{lemma:A_fst}
  \lean{A_fst}
  \leanok
  For any $c, b \in \mathbb{N}$ with $c \geq 4$, the first component of $A(c-4, b)$ satisfies:
  \[ A(c-4, b)_1 = C(c) - 4. \]
  \begin{proof}
    \leanok
    Let $a = c-4$. If $c$ is even, $a$ is even. Let $c=2k$, then $a=2k-4$ and $n = \lfloor a/2 \rfloor = k-2$.
    $A(a, b)_1 = 3(k-2)+2 = 3k-6+2 = 3k-4$.
    Also $C(c) = 3(2k)/2 = 3k$, so $C(c)-4 = 3k-4$.
    If $c$ is odd, $a$ is odd. Let $c=2k+1$, then $a=2k-3$ and $n = \lfloor a/2 \rfloor = k-2$.
    $A(a, b)_1 = 3(k-2)+3 = 3k-6+3 = 3k-3$.
    Also $C(c) = (3(2k+1)-1)/2 = (6k+2)/2 = 3k+1$, so $C(c)-4 = 3k-3$.
  \end{proof}
\end{lemma}

\begin{lemma}
  \label{lemma:C_ge_self}
  \lean{C_ge_self}
  \leanok
  For $n \geq 2$, we have $C(n) \geq n$.
  \begin{proof}
    \leanok
    $C(n) = \lfloor 3n/2 \rfloor$. Since $3n/2 \geq n$ for all $n \geq 0$, and $3n/2 \geq n+1$ is not always true,
    we check $n=2$: $C(2)=3 \geq 2$. $n=3$: $C(3)=4 \geq 3$.
    In general, $3n/2 - n = n/2$. For $n \geq 2$, $n/2 \geq 1$, so $3n/2 \geq n+1$ except for small $n$.
    The floor satisfies $\lfloor 3n/2 \rfloor > 3n/2 - 1 = n + n/2 - 1$.
    For $n \geq 2$, $n/2 - 1 \geq 0$, so $C(n) \geq n$.
  \end{proof}
\end{lemma}

\begin{lemma}
  \label{lemma:Citer_ge_8}
  \lean{Citer_ge_8}
  \leanok
  For all $i \in \mathbb{N}$, we have $C_i(8) \geq 8$.
  \begin{proof}
    \leanok
    By induction on $i$. Base case $i=0$: $C_0(8) = 8 \geq 8$.
    Inductive step: $C_{i+1}(8) = C(C_i(8))$. By induction $C_i(8) \geq 8$.
    Since $C$ is monotonic and $C(8) = 12 \geq 8$, the result follows from Lemma~\ref{lemma:C_ge_self}.
  \end{proof}
\end{lemma}

\begin{lemma}
  \label{lemma:Aiter_fst_eq}
  \lean{Aiter_fst_eq}
  \leanok
  For all $i \in \mathbb{N}$, the first component of $A_i(8, 2)$ satisfies:
  \[ A_i(8, 2)_1 = C_{i+1}(8) - 4. \]
  \begin{proof}
    \leanok
    By induction on $i$.
    Base case $i=0$: $A_0(8, 2)_1 = 8$. $C_1(8)-4 = C(8)-4 = 12-4 = 8$.
    Inductive step: $A_{i+1}(8, 2)_1 = A(A_i(8, 2))_1$.
    Since the first component of $A(a, b)$ depends only on $a$, and $A_i(8, 2)_1 = C_{i+1}(8)-4$,
    we apply Lemma~\ref{lemma:A_fst} with $c = C_{i+1}(8) \geq 8 \geq 4$:
    $A(C_{i+1}(8)-4, b)_1 = C(C_{i+1}(8)) - 4 = C_{i+2}(8) - 4$.
  \end{proof}
\end{lemma}

\begin{lemma}
  \label{lemma:Aiter_fst_parity}
  \lean{Aiter_fst_parity}
  \leanok
  For any $i \in \mathbb{N}$, $A_i(8, 2)_1$ is odd if and only if $C_{i+1}(8)$ is odd.
  \begin{proof}
    \leanok
    Follows from $A_i(8, 2)_1 = C_{i+1}(8) - 4$ and the fact that subtracting 4 preserves parity.
  \end{proof}
\end{lemma}

\begin{lemma}
  \label{lemma:Aiter_fst_div2}
  \lean{Aiter_fst_div2}
  \leanok
  For any $i \in \mathbb{N}$, $\lfloor A_i(8, 2)_1 / 2 \rfloor \geq 1$ if and only if $C_{i+1}(8) \geq 6$.
  \begin{proof}
    \leanok
    $A_i(8, 2)_1 / 2 \geq 1 \iff A_i(8, 2)_1 \geq 2 \iff C_{i+1}(8) - 4 \geq 2 \iff C_{i+1}(8) \geq 6$.
  \end{proof}
\end{lemma}

\begin{definition}
  \label{def:b_int}
  \lean{b_int}
  \leanok
  For $i \in \mathbb{N}$, we define:
  \[ b_{\text{int}}(i) = 2 + 2i - 3 \sum_{j=0}^i (C_{j+1}(8) \pmod 2). \]
\end{definition}

\begin{lemma}
  \label{lemma:Aiter_snd_eq_b_int}
  \lean{Aiter_snd_eq_b_int}
  \leanok
  Let $i \in \mathbb{N}$. If for all $j < i$, the condition $A_j(8, 2)_1 \equiv 1 \pmod 2$ implies $A_j(8, 2)_2 \geq 1$, then:
  \[ A_i(8, 2)_2 = b_{\text{int}}(i). \]
  \begin{proof}
    \leanok
    By induction on $i$. Base case $i=0$: $A_0(8, 2)_2 = 2$. $b_{\text{int}}(0) = 2 + 0 - 3(C_1(8) \pmod 2) = 2 - 3(0) = 2$.
    Step: $A_{i+1}(8, 2)_2 = A(A_i(8, 2))_2$.
    If $A_i(8, 2)_1$ is even, $A_{i+1}(8, 2)_2 = A_i(8, 2)_2 + 2$.
    If $A_i(8, 2)_1$ is odd, $A_{i+1}(8, 2)_2 = A_i(8, 2)_2 - 1$.
    In both cases, $A_{i+1}(8, 2)_2 = A_i(8, 2)_2 + 2 - 3(A_i(8, 2)_1 \pmod 2)$.
    Using parity equivalence and the induction hypothesis (and ensuring no truncation at 0), we get the formula.
  \end{proof}
\end{lemma}

\begin{lemma}
  \label{lemma:b_int_eq_zero_iff}
  \lean{b_int_eq_zero_iff}
  \leanok
  We have $b_{\text{int}}(i) = 0$ if and only if $i \equiv 2 \pmod 3$ and $\sum_{j=0}^i (C_{j+1}(8) \pmod 2) = \frac{2(i+1)}{3}$.
  \begin{proof}
    \leanok
    $b_{\text{int}}(i) = 0 \iff 2 + 2i = 3 \sum_{j=0}^i (C_{j+1}(8) \pmod 2)$.
    Since the RHS is a multiple of 3, we must have $2(i+1) \equiv 0 \pmod 3$, which means $i+1 \equiv 0 \pmod 3$, or $i \equiv 2 \pmod 3$.
    Then the sum must be $2(i+1)/3$.
  \end{proof}
\end{lemma}

\begin{theorem}
  \label{thm:Aiter_8_2_halt_iff_Citer_8_halt}
  \lean{Aiter_8_2_halt_iff_Citer_8_halt}
  \leanok
  The sequence $A_i(8, 2)$ halts at some index $i$ (meaning $A_i(8, 2)_1$ is odd and $A_i(8, 2)_2 = 0$) if and only if there exists $i$ such that $C_{i+1}(8)$ is odd, $i \equiv 2 \pmod 3$, and $\sum_{j=0}^i (C_{j+1}(8) \pmod 2) = \frac{2(i+1)}{3}$.
  \begin{proof}
    \leanok
    $\Rightarrow$: Let $i$ be the first index where $A_i(8, 2)_1$ is odd and $A_i(8, 2)_2 = 0$. Since it is the first such index, no truncation occurred before, so $A_i(8, 2)_2 = b_{\text{int}}(i) = 0$. The conditions then follow from Lemma~\ref{lemma:b_int_eq_zero_iff}.
    $\Leftarrow$: If the conditions hold for some $i$, then $b_{\text{int}}(i) = 0$. If no truncation occurs before $i$, then $A_i(8, 2)_2 = 0$ and we are done. If truncation occurred earlier, there was some $j < i$ where $A_j(8, 2)_2$ would have become negative but was clipped to 0, which also satisfies the halting condition.
  \end{proof}
\end{theorem}
